\begin{exsol}{30}{Cabral 13.10 Asset specificity and vertical integration}{346}{Advanced}{B3}
\asked Two firms $M$ and $R$ invest $e_M$, $e_R$ (in \$) in a joint project worth
$v = \sqrt{e_M} + \sqrt{e_R}$ and split $v$ equally after investing. Find (a) the equilibrium
investments, (b) the investments chosen by the merged firm, (c) total payoffs under separation
and integration, and (d) comment in light of \S13.1, Box~13.1 and \S5.1.

\method Simultaneous investment game with ex-post bargaining, solved by first-order conditions;
joint-surplus maximization for the merged firm. Hold-up and specific assets (Cabral \S13.1,
``Investment incentives'', pp.~338--339; \S3.4, p.~55); double marginalization and Box~13.1
(pp.~334--336); externalities (\S5.1, pp.~94--96). Course link: the hold-up preview in the
Aula~2 notes and \MR{} Ch.~5, investments and specific assets (pp.~134--139).

\solution
\subparagraph{(a) Vertical separation.} Investment is sunk before the split, and the split is
fixed at 50/50 whatever was invested. Firm $M$ solves
\[
\max_{e_M \ge 0}\ \tfrac12\big(\sqrt{e_M} + \sqrt{e_R}\big) - e_M .
\]
The objective is strictly concave in $e_M$; the first-order condition is
\[
\frac{1}{4\sqrt{e_M}} = 1 \;\Longrightarrow\; \sqrt{e_M} = \tfrac14 \;\Longrightarrow\;
e_M^{S} = \tfrac{1}{16}.
\]
$M$'s choice does not depend on $e_R$ (a dominant strategy), and by symmetry
$e_R^{S} = \tfrac1{16}$. Then $v^S = \tfrac14 + \tfrac14 = \tfrac12$, and each firm earns
$\tfrac12 \cdot \tfrac12 - \tfrac1{16} = \tfrac{3}{16}$.

\subparagraph{(b) Vertical integration.} Firm $M\&R$ keeps all of $v$ and pays both costs:
\[
\max_{e_M, e_R \ge 0}\ \sqrt{e_M} + \sqrt{e_R} - e_M - e_R
\quad\Longrightarrow\quad \frac{1}{2\sqrt{e_i}} = 1
\quad\Longrightarrow\quad e_M^{I} = e_R^{I} = \tfrac14 .
\]
Then $v^I = \tfrac12 + \tfrac12 = 1$. Each investment is four times the separated level.

\subparagraph{(c) Total payoffs.}
\[
W^{S} = v^S - e_M^S - e_R^S = \tfrac12 - \tfrac18 = \tfrac38, \qquad
W^{I} = v^I - e_M^I - e_R^I = 1 - \tfrac12 = \tfrac12 .
\]
Integration raises total payoff by $\tfrac12 - \tfrac38 = \tfrac18$, a gain of one third over
separation (a quarter of the efficient surplus is lost under separation).
\begin{center}\small
\begin{tabular}{@{}lcccc@{}}
\toprule
& $e_M$ & $e_R$ & $v$ & Total payoff \\
\midrule
Separation, 50/50 split & $1/16$ & $1/16$ & $1/2$ & $3/8$ \\
Integration & $1/4$ & $1/4$ & $1$ & $1/2$ \\
\bottomrule
\end{tabular}
\end{center}

\subparagraph{(d) Comment.}
\emph{Hold-up (\S13.1).} Each investment is specific: it creates value only inside the
relationship, and once sunk it no longer affects how $v$ is divided. Each firm pays the full
marginal cost of investing but receives only half the marginal value, so it equates
$\tfrac12 \cdot \tfrac{1}{2\sqrt{e_i}}$ with $1$ instead of $\tfrac{1}{2\sqrt{e_i}}$ with $1$.
This is the continuous version of Cabral's hold-up example (p.~338): there, an investment worth
$8$ and costing $5$ was not made because the investor would get only $4$ in an even split
(p.~339); here investment is made, but at a quarter of the efficient level. The margin note on
p.~339 says this exercise is ``the case of a continuum of choices''. Integration puts both
decisions under one owner who receives the whole return, which removes the distortion; Cabral
adds that it creates an agency problem in its place, since the acquired unit's manager now
follows orders with weaker incentives (pp.~338--339). Integration is worth it only if that
agency cost is below the gain $\tfrac18$.

\emph{Externalities (\S5.1).} The 50/50 split makes each firm's investment confer a positive
externality on the other: half of the extra value goes to the partner, who did not pay for it.
As with any positive externality (front yards, item~9), the activity is under-provided. The
Coase remedy (bargain to the efficient level) would require contracting on $e_M$ and $e_R$
\emph{before} they are sunk, which is exactly what is assumed impossible: investments are hard
to specify and verify, so contracts are incomplete (p.~338). Ex-post bargaining is efficient
given the investments; it is the ex-ante investments that are distorted.

\emph{Box~13.1 and double marginalization.} Box~13.1 (p.~336) reports that integrated
distributors and sports networks charge consumers lower prices, evidence that integration
removes the double-marginalization externality: under separation each firm adds its own margin
and ignores the effect on the other's profit (pp.~334--336). The hold-up problem is the
investment counterpart of the same failure: separated firms ignore the effect of their choice on
the partner's payoff, and integration internalizes it in both cases.

\interp The exercise isolates why ownership matters even when bargaining after the fact is
efficient: the anticipated split determines the incentive to invest beforehand. This is the core
of the transaction-cost and property-rights theories of the firm and the Topic~6 material that
the Aula~2 notes preview (their quadratic example gives the same $\tfrac38$ against
$\tfrac12$). It also explains the Smart-car strain of item~13: a jointly owned venture with
investments by both partners is a 50/50 split.

\verify (1) \texttt{check\_B.py} confirms with sympy: $e_i^S = \tfrac1{16}$, payoff
$\tfrac3{16}$ each, $W^S = \tfrac38$; $e_i^I = \tfrac14$, $W^I = \tfrac12$; difference
$\tfrac18$, a one-third gain. (2) \emph{No split fixes it.} If $M$ receives a share $s$ and $R$
receives $1-s$, then $e_M = s^2/4$, $e_R = (1-s)^2/4$ and
$W(s) = \tfrac12 - \tfrac14\big[s^2 + (1-s)^2\big]$, maximized at $s=\tfrac12$ with
$W = \tfrac38 < \tfrac12$; giving everything to one firm ($s=1$) yields only $\tfrac14$, because
the other then stops investing. With two-sided investment, no ex-post division restores
efficiency; joint ownership of both decisions does. (3) \emph{One-sided case:} if only $M$
invested, giving $M$ the whole value ($s=1$) would induce the efficient $e_M = \tfrac14$, so
ownership should go to the investor. That is the property-rights logic behind the exercise.

\answer{(a) $e_M^S = e_R^S = 1/16$ (from $\tfrac{1}{4\sqrt{e_i}} = 1$), payoff $3/16$ each.
(b) $e_M^I = e_R^I = 1/4$. (c) $W^S = 3/8 < W^I = 1/2$: integration gains $1/8$ (one third).
(d) A 50/50 ex-post split gives each investor half the marginal value of its sunk, specific
investment, a positive externality on the partner and a continuous hold-up problem (\S13.1,
\S5.1). Integration internalizes it, as it internalizes double marginalization (Box~13.1),
at the cost of new agency problems.}
\end{exsol}
